\documentclass[11pt,a4paper]{article}
\usepackage{fontspec,amsmath,amssymb,amsthm,microtype}
\setmainfont{lmroman10-regular.otf}[BoldFont=lmroman10-bold.otf,ItalicFont=lmroman10-italic.otf,BoldItalicFont=lmroman10-bolditalic.otf]
\usepackage{xcolor}
\usepackage[most]{tcolorbox}
\usepackage{mathrsfs}
\providecommand{\lt}{<}
\providecommand{\gt}{>}
\setmonofont{DejaVuSansMono.ttf}[BoldFont=DejaVuSansMono-Bold.ttf,ItalicFont=DejaVuSansMono-Oblique.ttf,BoldItalicFont=DejaVuSansMono-BoldOblique.ttf,Scale=0.86]
\newfontfamily\leanfallback{STIXTwoMath-Regular.otf}[Scale=0.86]
\definecolor{leanrule}{gray}{0.62}
\definecolor{leanback}{gray}{0.962}
\definecolor{leantitle}{gray}{0.22}
\definecolor{leanmuted}{gray}{0.38}
\newtcolorbox{leanbox}[1]{enhanced,breakable,colback=leanback,colframe=leanrule,
  boxrule=0.4pt,arc=0.6mm,left=2.4mm,right=2.4mm,top=1.8mm,bottom=1.8mm,
  before skip=0.8em,after skip=1em,fonttitle=\small\bfseries,coltitle=leantitle,
  colbacktitle=leanback,titlerule=0.3pt,title={#1}}
\newlength{\leancharwidth}
\newenvironment{leancode}{\par\smallskip\begingroup\ttfamily\footnotesize\raggedright
  \settowidth{\leancharwidth}{M}\setlength{\parskip}{0pt}\setlength{\parindent}{0pt}
  \frenchspacing}{\par\endgroup\smallskip}
\newcommand{\leanline}[2]{\hangindent=\dimexpr#1\leancharwidth+4\leancharwidth\relax
  \hangafter=1\noindent\hspace*{#1\leancharwidth}#2\par}
\newcommand{\leansource}[1]{\par{\footnotesize\color{leanmuted}#1}\par}
\newcommand{\leanlabel}[1]{\par\smallskip\noindent{\small\bfseries #1}\par}
\usepackage[margin=25mm]{geometry}
\usepackage[unicode,hidelinks]{hyperref}
\hypersetup{pdfauthor={LangevinLean Project}}
\setlength{\emergencystretch}{3em}
\setlength{\parskip}{0.4em}
\setlength{\parindent}{0pt}
\allowdisplaybreaks[2]
\raggedbottom
\title{\texorpdfstring{The gradient map on square-integrable random vectors}{The gradient map on square-integrable random vectors}}
\author{LangevinLean Project}
\date{}
\begin{document}\maketitle
Draft mathematical exposition; these notes may contain mistakes.


Let \(P\) be a finite measure on a measurable space \(\Omega\) and \(H\) a normed space (from Step 6 on, a real Hilbert space). We write \(L^2(P;H)\) for the space of equivalence classes, modulo equality \(P\)-almost everywhere, of almost everywhere strongly measurable functions \(X:\Omega\to H\) with \(\int\|X\|^2\,dP<\infty\), normed by \(\|X\|_{L^2}=(\int\|X\|^2\,dP)^{1/2}\). For a class \(X\) we write \(X(\omega)\) for the values of a fixed representative; every identity between such values holds for \(P\)-almost every \(\omega\) and does not depend on the representative. When \(H\) is a Hilbert space, so is \(L^2(P;H)\), with inner product \(\langle X,Y\rangle=\int\langle X(\omega),Y(\omega)\rangle\,dP(\omega)\).




For an \(L\)-Lipschitz map \(g:H\to H\) we construct the map \(\mathcal G:L^2(P;H)\to L^2(P;H)\) with \(\mathcal G(X)(\omega)=g(X(\omega))\) almost everywhere (the lifted drift) and prove that it is again \(L\)-Lipschitz and that \(\|\mathcal G(X)\|_{L^2}^2=\int\|g(X)\|^2\,dP\). Finiteness of \(P\) is needed because \(g(0)\) need not vanish and constant functions must be square integrable. For \(g=\nabla U\), with \(U\) strongly convex with Lipschitz gradient, the sharp contraction factors \(1-mh\) and \(Lh-1\) of the gradient step carry over to \(L^2\), and the local-defect bound applies in the Banach space \(L^2(P;H)\), where time integrals are Bochner integrals of \(L^2\)-valued curves. Finally the ULA update \(X-h\mathcal G(X)+\sqrt{2h}\,Z\) and the resulting chain are defined on \(L^2(P;E_d)\), and every iterate is shown to agree almost everywhere with the pointwise ULA recursion driven by the representatives of the initial state and the noise. In this way the coupling argument for random vectors becomes an argument in a Hilbert space whose norm is the root-mean-square distance.


\section{\texorpdfstring{Centring a Lipschitz map}{Centring a Lipschitz map}}


\textbf{Claim.} Let \(H\) be a normed additive group, \(L\ge0\), and \(g:H\to H\) be \(L\)-Lipschitz. Then \(\tilde g(x)=g(x)-g(0)\) is \(L\)-Lipschitz.




\textbf{Proof.} For \(x,y\in H\), \(\tilde g(x)-\tilde g(y)=(g(x)-g(0))-(g(y)-g(0))=g(x)-g(y)\), so \(\|\tilde g(x)-\tilde g(y)\|=\|g(x)-g(y)\|\le L\|x-y\|\). \(\square\)




Moreover \(\tilde g(0)=0\), so \(\|\tilde g(x)\|=\|\tilde g(x)-\tilde g(0)\|\le L\|x\|\) for every \(x\).

\begin{leanbox}{Lean: \texttt{centered\_\allowbreak{}lipschitz}}
\leansource{\texttt{centered\_\allowbreak{}lipschitz} \quad UnadjustedLangevin/\allowbreak{}LpDrift.lean, lines 26–30}
\begin{leancode}
\leanline{0}{theorem\ centered\_lipschitz\ \{L\ :\ ℝ≥0\}\ \{g\ :\ H\ →\ H\}\ (hLip\ :\ LipschitzWith\ L\ g)\ :}
\leanline{4}{LipschitzWith\ L\ (fun\ x\ ↦\ g\ x\ -\ g\ 0)}
\end{leancode}
\leanlabel{How the Lean proof follows the argument}
\texttt{centered\_\allowbreak{}lipschitz} uses \texttt{Lipschitz\allowbreak{}With.\allowbreak{}of\_\allowbreak{}dist\_\allowbreak{}le\_\allowbreak{}mul} (the Lipschitz property in the form \(\operatorname{dist}(\tilde g(x),\tilde g(y))\le L\operatorname{dist}(x,y)\)) and closes the resulting inequality with \texttt{h\allowbreak{}Lip.\allowbreak{}dist\_\allowbreak{}le\_\allowbreak{}mul x y}, after rewriting distances as norms (\texttt{dist\_\allowbreak{}eq\_\allowbreak{}norm}) and cancelling the offset (\texttt{sub\_\allowbreak{}sub\_\allowbreak{}sub\_\allowbreak{}cancel\_\allowbreak{}right}, i.e. \((a-c)-(b-c)=a-b\)).
\end{leanbox}

\section{\texorpdfstring{The lifted drift}{The lifted drift}}


\textbf{Definition.} Let \(P\) be a finite measure, \(L\ge0\) and \(g:H\to H\) be \(L\)-Lipschitz. For \(X\in L^2(P;H)\) put


\[\mathcal G(X)=[\tilde g\circ X]+[g(0)],\]


where \(\tilde g=g-g(0)\) and \([g(0)]\) is the class of the constant function.




\emph{Both terms lie in \(L^2(P;H)\).} The function \(\tilde g\) is continuous, so \(\tilde g\circ X\) is almost everywhere strongly measurable, and by the previous step \(\|\tilde g(X(\omega))\|^2\le L^2\|X(\omega)\|^2\), hence \(\int\|\tilde g\circ X\|^2\,dP\le L^2\|X\|_{L^2}^2<\infty\). If two representatives of \(X\) agree almost everywhere, so do their compositions with \(\tilde g\); thus \([\tilde g\circ X]\) depends only on the class. The constant function has \(\int\|g(0)\|^2\,dP=\|g(0)\|^2P(\Omega)<\infty\) because \(P\) is finite.

\begin{leanbox}{Lean: \texttt{lpDrift}}
\leansource{\texttt{lpDrift} \quad UnadjustedLangevin/\allowbreak{}LpDrift.lean, lines 34–36}
\begin{leancode}
\leanline{0}{def\ lpDrift\ (g\ :\ H\ →\ H)\ \{L\ :\ ℝ≥0\}\ (hLip\ :\ LipschitzWith\ L\ g)}
\leanline{4}{(X\ :\ Lp\ H\ 2\ P)\ :\ Lp\ H\ 2\ P\ :=}
\leanline{2}{(centered\_lipschitz\ hLip).compLp\ (by\ simp)\ X\ +\ Lp.const\ 2\ P\ (g\ 0)}
\end{leancode}
\leanlabel{How the Lean proof follows the argument}
\texttt{lp\allowbreak{}Drift P g h\allowbreak{}Lip X} is defined as \texttt{(centered\_\allowbreak{}lipschitz h\allowbreak{}Lip).\allowbreak{}comp\allowbreak{}Lp (by simp) X + Lp.\allowbreak{}const 2 P (g 0)}. Here \texttt{Lp H 2 P} is Mathlib's space \(L^2(P;H)\) of almost-everywhere classes. Mathlib's \texttt{Lipschitz\allowbreak{}With.\allowbreak{}comp\allowbreak{}Lp} composes an element of \texttt{Lp} with a Lipschitz map vanishing at \(0\); its side condition \(\tilde g(0)=0\) is proved by \texttt{simp}. \texttt{Lp.\allowbreak{}const 2 P (g 0)} is the class of the constant function, available because of the instance \texttt{[Is\allowbreak{}Finite\allowbreak{}Measure P]}.
\end{leanbox}

\section{\texorpdfstring{The lift acts pointwise}{The lift acts pointwise}}


\textbf{Claim.} For a finite measure \(P\), an \(L\)-Lipschitz \(g\) and \(X\in L^2(P;H)\),


\[\mathcal G(X)(\omega)=g(X(\omega))\qquad\text{for }P\text{-almost every }\omega.\]


\textbf{Proof.} Each of the following holds outside a \(P\)-null set: the representative of the sum \(\mathcal G(X)\) is the sum of the representatives of the two terms; the representative of \([\tilde g\circ X]\) equals \(\tilde g(X(\omega))\); the representative of \([g(0)]\) equals \(g(0)\). A union of three null sets is null, and outside it


\[\mathcal G(X)(\omega)=\bigl(g(X(\omega))-g(0)\bigr)+g(0)=g(X(\omega)).\qquad\square\]\begin{leanbox}{Lean: \texttt{lpDrift\_\allowbreak{}coe}}
\leansource{\texttt{lpDrift\_\allowbreak{}coe} \quad UnadjustedLangevin/\allowbreak{}LpDrift.lean, lines 39–47}
\begin{leancode}
\leanline{0}{theorem\ lpDrift\_coe\ (g\ :\ H\ →\ H)\ \{L\ :\ ℝ≥0\}\ (hLip\ :\ LipschitzWith\ L\ g)}
\leanline{4}{(X\ :\ Lp\ H\ 2\ P)\ :\ lpDrift\ P\ g\ hLip\ X\ =ᵐ[P]\ fun\ ω\ ↦\ g\ (X\ ω)}
\end{leancode}
\leanlabel{How the Lean proof follows the argument}
In \texttt{lp\allowbreak{}Drift\_\allowbreak{}coe}, \texttt{=ᵐ[P]} denotes equality \(P\)-almost everywhere. After unfolding, \texttt{filter\_\allowbreak{}upwards} intersects the three almost-everywhere statements \texttt{Lp.\allowbreak{}coe\allowbreak{}Fn\_\allowbreak{}add}, \texttt{coe\allowbreak{}Fn\_\allowbreak{}comp\allowbreak{}Lp} and \texttt{Lp.\allowbreak{}coe\allowbreak{}Fn\_\allowbreak{}const} (the representatives of a sum, of a composition and of a constant), fixes \(\omega\) in the intersection, and \texttt{sub\_\allowbreak{}add\_\allowbreak{}cancel} cancels the offset.
\end{leanbox}

\section{\texorpdfstring{Square integrability of the drift}{Square integrability of the drift}}


\textbf{Claim.} For a finite measure \(P\), an \(L\)-Lipschitz \(g\) and \(X\in L^2(P;H)\), the function \(\omega\mapsto g(X(\omega))\) is square integrable.




\textbf{Proof.} The class \(\mathcal G(X)\) lies in \(L^2(P;H)\), so its representative is square integrable. By the previous step it agrees almost everywhere with \(g\circ X\), and square integrability is preserved under almost-everywhere equality. \(\square\)

\begin{leanbox}{Lean: \texttt{lpDrift\_\allowbreak{}memLp}}
\leansource{\texttt{lpDrift\_\allowbreak{}memLp} \quad UnadjustedLangevin/\allowbreak{}LpDrift.lean, lines 50–52}
\begin{leancode}
\leanline{0}{theorem\ lpDrift\_memLp\ (g\ :\ H\ →\ H)\ \{L\ :\ ℝ≥0\}\ (hLip\ :\ LipschitzWith\ L\ g)}
\leanline{4}{(X\ :\ Lp\ H\ 2\ P)\ :\ MemLp\ (fun\ ω\ ↦\ g\ (X\ ω))\ 2\ P}
\end{leancode}
\leanlabel{How the Lean proof follows the argument}
\texttt{lp\allowbreak{}Drift\_\allowbreak{}mem\allowbreak{}Lp} is the term \texttt{(Lp.\allowbreak{}mem\allowbreak{}Lp (lp\allowbreak{}Drift P g h\allowbreak{}Lip X)).\allowbreak{}ae\_\allowbreak{}eq (lp\allowbreak{}Drift\_\allowbreak{}coe P g h\allowbreak{}Lip X)}: \texttt{Lp.\allowbreak{}mem\allowbreak{}Lp} gives square integrability of the representative and \texttt{Mem\allowbreak{}Lp.\allowbreak{}ae\_\allowbreak{}eq} transfers it along the almost-everywhere equality.
\end{leanbox}

\section{\texorpdfstring{The lifted drift is Lipschitz}{The lifted drift is Lipschitz}}


\textbf{Claim.} For a finite measure \(P\) and an \(L\)-Lipschitz \(g\), the map \(\mathcal G\) is \(L\)-Lipschitz on \(L^2(P;H)\):


\[\|\mathcal G(X)-\mathcal G(Y)\|_{L^2}\le L\,\|X-Y\|_{L^2}\qquad(X,Y\in L^2(P;H)).\]


\textbf{Proof.} The constant terms cancel: \(\mathcal G(X)-\mathcal G(Y)=[\tilde g\circ X]-[\tilde g\circ Y]\). For almost every \(\omega\), the representative of this difference is \(\tilde g(X(\omega))-\tilde g(Y(\omega))\) and that of \(X-Y\) is \(X(\omega)-Y(\omega)\), and by Step 1


\[\|\tilde g(X(\omega))-\tilde g(Y(\omega))\|^2\le L^2\|X(\omega)-Y(\omega)\|^2.\]


Integrating over \(\Omega\) and taking square roots gives the claim. \(\square\)




The constant \(g(0)\) does not contribute to the difference, so no term proportional to \(\|g(0)\|\) appears.

\begin{leanbox}{Lean: \texttt{lpDrift\_\allowbreak{}lipschitz}}
\leansource{\texttt{lpDrift\_\allowbreak{}lipschitz} \quad UnadjustedLangevin/\allowbreak{}LpDrift.lean, lines 55–60}
\begin{leancode}
\leanline{0}{theorem\ lpDrift\_lipschitz\ (g\ :\ H\ →\ H)\ \{L\ :\ ℝ≥0\}\ (hLip\ :\ LipschitzWith\ L\ g)\ :}
\leanline{4}{LipschitzWith\ L\ (lpDrift\ P\ g\ hLip)}
\end{leancode}
\leanlabel{How the Lean proof follows the argument}
\texttt{lp\allowbreak{}Drift\_\allowbreak{}lipschitz} reduces \texttt{Lipschitz\allowbreak{}With L (lp\allowbreak{}Drift P g h\allowbreak{}Lip)} to the distance inequality (\texttt{Lipschitz\allowbreak{}With.\allowbreak{}of\_\allowbreak{}dist\_\allowbreak{}le\_\allowbreak{}mul}), cancels the constants by \texttt{add\_\allowbreak{}sub\_\allowbreak{}add\_\allowbreak{}right\_\allowbreak{}eq\_\allowbreak{}sub}, and applies Mathlib's \texttt{norm\_\allowbreak{}comp\allowbreak{}Lp\_\allowbreak{}sub\_\allowbreak{}le} for \texttt{comp\allowbreak{}Lp}, which is the integrated pointwise estimate above.
\end{leanbox}

\section{\texorpdfstring{The small-step contraction in \(L^2\)}{The small-step contraction in L²}}


\textbf{Claim.} Let \(H\) be a real Hilbert space, \(P\) a finite measure, \(m,h\in\mathbb R\), \(L\ge0\), and \(U:H\to\mathbb R\) differentiable and \(m\)-strongly convex with \(L\)-Lipschitz gradient. Assume \(m\le L\), \(m+L>0\) and \(0\le h\le2/(m+L)\). Let \(\mathcal G\) be the lift of \(g=\nabla U\). Then for all \(X,Y\in L^2(P;H)\),


\[\bigl\|(X-h\mathcal G(X))-(Y-h\mathcal G(Y))\bigr\|_{L^2}\le(1-mh)\,\|X-Y\|_{L^2}.\]


\textbf{Proof.} \emph{(a) Pointwise bounds integrate.} Let \(F,G\in L^2(P;H)\) and \(c\in\mathbb R\) with \(\|F(\omega)\|\le c\,\|G(\omega)\|\) for almost every \(\omega\). If \(c\ge0\), squaring and integrating gives \(\int\|F\|^2\le c^2\int\|G\|^2\), hence \(\|F\|_{L^2}\le c\|G\|_{L^2}\). If \(c<0\), the inequality forces \(F(\omega)=0\) and \(G(\omega)=0\) almost everywhere, and both sides vanish.




\emph{(b) The pointwise bound.} Apply (a) with \(F=(X-h\mathcal G(X))-(Y-h\mathcal G(Y))\), \(G=X-Y\) and \(c=1-mh\). Outside a null set (the union of finitely many null sets on which the representatives of differences and scalar multiples, and the formula of Step 3 for \(\mathcal G(X)\) and \(\mathcal G(Y)\), may fail),


\[F(\omega)=\bigl(X(\omega)-h\nabla U(X(\omega))\bigr)-\bigl(Y(\omega)-h\nabla U(Y(\omega))\bigr),\qquad G(\omega)=X(\omega)-Y(\omega).\]


The sharp small-step contraction of the gradient step at \(x=X(\omega)\), \(y=Y(\omega)\) gives \(\|F(\omega)\|\le(1-mh)\|G(\omega)\|\). \(\square\)




The contraction acts on random states through their representatives; no assumption on the joint law of \(X\) and \(Y\) is involved.

\begin{leanbox}{Lean: \texttt{lpDrift\_\allowbreak{}gradient\_\allowbreak{}contract}}
\leansource{\texttt{lpDrift\_\allowbreak{}gradient\_\allowbreak{}contract} \quad UnadjustedLangevin/\allowbreak{}LpDrift.lean, lines 66–86}
\begin{leancode}
\leanline{0}{theorem\ lpDrift\_gradient\_contract}
\leanline{4}{(U\ :\ H\ →\ ℝ)\ (m\ h\ :\ ℝ)\ (L\ :\ ℝ≥0)}
\leanline{4}{(hdiff\ :\ Differentiable\ ℝ\ U)\ (hsc\ :\ StrongConvexOn\ Set.univ\ m\ U)}
\leanline{4}{(hLip\ :\ LipschitzWith\ L\ (gradient\ U))\ (hmL\ :\ m\ ≤\ (L\ :\ ℝ))}
\leanline{4}{(hpos\ :\ 0\ <\ m\ +\ (L\ :\ ℝ))\ (hh\ :\ 0\ ≤\ h)\ (hsmall\ :\ h\ ≤\ 2\ /\ (m\ +\ (L\ :\ ℝ)))}
\leanline{4}{(X\ Y\ :\ Lp\ H\ 2\ P)\ :}
\leanline{4}{‖(X\ -\ h\ •\ lpDrift\ P\ (gradient\ U)\ hLip\ X)\ -}
\leanline{6}{(Y\ -\ h\ •\ lpDrift\ P\ (gradient\ U)\ hLip\ Y)‖\ ≤\ (1\ -\ m\ *\ h)\ *\ ‖X\ -\ Y‖}
\end{leancode}
\leanlabel{How the Lean proof follows the argument}
\texttt{lp\allowbreak{}Drift\_\allowbreak{}gradient\_\allowbreak{}contract} applies Mathlib's \texttt{Lp.\allowbreak{}norm\_\allowbreak{}le\_\allowbreak{}mul\_\allowbreak{}norm\_\allowbreak{}of\_\allowbreak{}ae\_\allowbreak{}le\_\allowbreak{}mul}, which is (a) for an arbitrary real constant. \texttt{filter\_\allowbreak{}upwards} intersects eight almost-everywhere statements: \texttt{Lp.\allowbreak{}coe\allowbreak{}Fn\_\allowbreak{}sub} for the outer difference, for both Euler differences and for \texttt{X -\allowbreak{} Y}, \texttt{Lp.\allowbreak{}coe\allowbreak{}Fn\_\allowbreak{}smul} for both scalar multiples, and \texttt{lp\allowbreak{}Drift\_\allowbreak{}coe} for \texttt{X} and \texttt{Y}. After \texttt{simp only} substitutes these representatives, the goal is closed by \texttt{strong\_\allowbreak{}smooth\_\allowbreak{}gradient\_\allowbreak{}step\_\allowbreak{}contract} at the points \texttt{X \ensuremath{\omega}} and \texttt{Y \ensuremath{\omega}}.
\end{leanbox}

\section{\texorpdfstring{The large-step bound in \(L^2\)}{The large-step bound in L²}}


\textbf{Claim.} With \(H\), \(P\), \(U\), \(m\), \(L\) and \(\mathcal G\) as in the previous step, assume \(m\le L\), \(m+L>0\), \(h\ge0\) and \(2/(m+L)\le h\). Then for all \(X,Y\in L^2(P;H)\),


\[\bigl\|(X-h\mathcal G(X))-(Y-h\mathcal G(Y))\bigr\|_{L^2}\le(Lh-1)\,\|X-Y\|_{L^2}.\]


\textbf{Proof.} As in the previous step, with the constant \(c=Lh-1\) and the sharp large-step bound of the gradient step in place of the small-step bound. \(\square\)

\begin{leanbox}{Lean: \texttt{lpDrift\_\allowbreak{}gradient\_\allowbreak{}contract\_\allowbreak{}large}}
\leansource{\texttt{lpDrift\_\allowbreak{}gradient\_\allowbreak{}contract\_\allowbreak{}large} \quad UnadjustedLangevin/\allowbreak{}LpDrift.lean, lines 89–109}
\begin{leancode}
\leanline{0}{theorem\ lpDrift\_gradient\_contract\_large}
\leanline{4}{(U\ :\ H\ →\ ℝ)\ (m\ h\ :\ ℝ)\ (L\ :\ ℝ≥0)}
\leanline{4}{(hdiff\ :\ Differentiable\ ℝ\ U)\ (hsc\ :\ StrongConvexOn\ Set.univ\ m\ U)}
\leanline{4}{(hLip\ :\ LipschitzWith\ L\ (gradient\ U))\ (hmL\ :\ m\ ≤\ (L\ :\ ℝ))}
\leanline{4}{(hpos\ :\ 0\ <\ m\ +\ (L\ :\ ℝ))\ (hh\ :\ 0\ ≤\ h)\ (hlarge\ :\ 2\ /\ (m\ +\ (L\ :\ ℝ))\ ≤\ h)}
\leanline{4}{(X\ Y\ :\ Lp\ H\ 2\ P)\ :}
\leanline{4}{‖(X\ -\ h\ •\ lpDrift\ P\ (gradient\ U)\ hLip\ X)\ -}
\leanline{6}{(Y\ -\ h\ •\ lpDrift\ P\ (gradient\ U)\ hLip\ Y)‖\ ≤\ ((L\ :\ ℝ)\ *\ h\ -\ 1)\ *\ ‖X\ -\ Y‖}
\end{leancode}
\leanlabel{How the Lean proof follows the argument}
\texttt{lp\allowbreak{}Drift\_\allowbreak{}gradient\_\allowbreak{}contract\_\allowbreak{}large} has the same proof as the previous step, ending with \texttt{strong\_\allowbreak{}smooth\_\allowbreak{}gradient\_\allowbreak{}step\_\allowbreak{}contract\_\allowbreak{}large} at \texttt{X \ensuremath{\omega}} and \texttt{Y \ensuremath{\omega}}.
\end{leanbox}

\section{\texorpdfstring{The local defect in \(L^2\)}{The local defect in L²}}


\textbf{Claim.} Let \(H\) be a real Hilbert space, \(P\) a finite measure, \(L\ge0\), \(U:H\to\mathbb R\) with \(L\)-Lipschitz gradient, and \(\mathcal G\) the lift of \(\nabla U\). Let \(Y,B:\mathbb R\to L^2(P;H)\) and \(h,d\in\mathbb R\) with \(h\ge0\), \(d\ge0\) and \(Lh\le2\). Assume that for every \(t\in[0,h]\)


\[\|\mathcal G(Y_t)\|_{L^2}\le\sqrt{Ld},\qquad\|B_t\|_{L^2}\le\sqrt{2d}\,\sqrt t,\qquad Y_t-Y_0=-\int_0^t\mathcal G(Y_s)\,ds+B_t,\]


the integral being the Bochner integral of an \(L^2(P;H)\)-valued function. Then


\[\|D_h(\mathcal G,Y)\|_{L^2}\le\frac{33}{20}\,Lh\sqrt{hd},\qquad D_h(\mathcal G,Y)=\int_0^h\bigl(\mathcal G(Y_t)-\mathcal G(Y_0)\bigr)\,dt.\]


\textbf{Proof.} The space \(L^2(P;H)\) is a real Banach space, and by Step 5 the map \(\mathcal G\) is \(L\)-Lipschitz on it. The hypotheses are those of the local-defect bound with the stationary moment scales, applied in the Banach space \(L^2(P;H)\) to the drift \(\mathcal G\); that bound gives the claim. \(\square\)




Differentiability and convexity of \(U\) are not needed here; only the Lipschitz constant of the gradient enters. The norms are root-mean-square norms: \(\|\mathcal G(Y_t)\|_{L^2}^2=\int\|\nabla U(Y_t)\|^2dP\) by the next step.

\begin{leanbox}{Lean: \texttt{lpDrift\_\allowbreak{}localDefect\_\allowbreak{}bound}}
\leansource{\texttt{lpDrift\_\allowbreak{}localDefect\_\allowbreak{}bound} \quad UnadjustedLangevin/\allowbreak{}LpDrift.lean, lines 112–124}
\begin{leancode}
\leanline{0}{theorem\ lpDrift\_localDefect\_bound}
\leanline{4}{(U\ :\ H\ →\ ℝ)\ (L\ :\ ℝ≥0)\ (h\ d\ :\ ℝ)}
\leanline{4}{(hLip\ :\ LipschitzWith\ L\ (gradient\ U))\ (Y\ B\ :\ ℝ\ →\ Lp\ H\ 2\ P)}
\leanline{4}{(hh\ :\ 0\ ≤\ h)\ (hd\ :\ 0\ ≤\ d)\ (hstep\ :\ (L\ :\ ℝ)\ *\ h\ ≤\ 2)}
\leanline{4}{(hG\ :\ ∀\ t\ ∈\ Set.Icc\ 0\ h,\ ‖lpDrift\ P\ (gradient\ U)\ hLip\ (Y\ t)‖\ ≤}
\leanline{6}{Real.sqrt\ ((L\ :\ ℝ)\ *\ d))}
\leanline{4}{(hB\ :\ ∀\ t\ ∈\ Set.Icc\ 0\ h,\ ‖B\ t‖\ ≤\ Real.sqrt\ (2\ *\ d)\ *\ Real.sqrt\ t)}
\leanline{4}{(heq\ :\ ∀\ t\ ∈\ Set.Icc\ 0\ h,\ Y\ t\ -\ Y\ 0\ =}
\leanline{6}{-(∫\ s\ in\ 0..t,\ lpDrift\ P\ (gradient\ U)\ hLip\ (Y\ s))\ +\ B\ t)\ :}
\leanline{4}{‖localDefect\ (lpDrift\ P\ (gradient\ U)\ hLip)\ Y\ h‖\ ≤}
\leanline{6}{(33\ /\ 20\ :\ ℝ)\ *\ (L\ :\ ℝ)\ *\ h\ *\ Real.sqrt\ (h\ *\ d)}
\end{leancode}
\leanlabel{How the Lean proof follows the argument}
\texttt{lp\allowbreak{}Drift\_\allowbreak{}local\allowbreak{}Defect\_\allowbreak{}bound} is \texttt{local\allowbreak{}Defect\_\allowbreak{}dalalyan\_\allowbreak{}bound} instantiated with the drift \texttt{lp\allowbreak{}Drift P (gradient U) h\allowbreak{}Lip} on the space \texttt{Lp H 2 P}, whose Lipschitz property is \texttt{lp\allowbreak{}Drift\_\allowbreak{}lipschitz P (gradient U) h\allowbreak{}Lip}. The integral \texttt{\ensuremath{\int} s in 0.\allowbreak{}.\allowbreak{}t, .\allowbreak{}.\allowbreak{}.\allowbreak{}} in \texttt{heq} is a Bochner integral in \texttt{Lp H 2 P}, which is complete.
\end{leanbox}

\section{\texorpdfstring{The drift norm as a second moment}{The drift norm as a second moment}}


\textbf{Claim.} Let \(H\) be a real inner-product space, \(P\) a finite measure, \(g:H\to H\) \(L\)-Lipschitz and \(X\in L^2(P;H)\). Then


\[\|\mathcal G(X)\|_{L^2}^2=\int_\Omega\|g(X(\omega))\|^2\,P(d\omega).\]


\textbf{Proof.} By definition of the inner product of \(L^2\),

\[\begin{aligned}\|\mathcal G(X)\|_{L^2}^2&=\langle\mathcal G(X),\mathcal G(X)\rangle\\&=\int\langle\mathcal G(X)(\omega),\mathcal G(X)(\omega)\rangle\,dP.\end{aligned}\]

By Step 3 the integrand equals \(\langle g(X(\omega)),g(X(\omega))\rangle=\|g(X(\omega))\|^2\) almost everywhere, and integrals of almost-everywhere equal functions coincide. \(\square\)




If \(X\) has law \(\pi\) under \(P\), the right-hand side equals \(\int\|g(x)\|^2\,\pi(dx)\); for \(g=\nabla U\) and \(\pi\) the Gibbs measure this is the score moment.

\begin{leanbox}{Lean: \texttt{lpDrift\_\allowbreak{}norm\_\allowbreak{}sq}}
\leansource{\texttt{lpDrift\_\allowbreak{}norm\_\allowbreak{}sq} \quad UnadjustedLangevin/\allowbreak{}LpDrift.lean, lines 128–134}
\begin{leancode}
\leanline{0}{theorem\ lpDrift\_norm\_sq\ (g\ :\ H\ →\ H)\ \{L\ :\ ℝ≥0\}\ (hLip\ :\ LipschitzWith\ L\ g)}
\leanline{4}{(X\ :\ Lp\ H\ 2\ P)\ :}
\leanline{4}{‖lpDrift\ P\ g\ hLip\ X‖\ \textasciicircum{}\ 2\ =\ ∫\ ω,\ ‖g\ (X\ ω)‖\ \textasciicircum{}\ 2\ ∂P}
\end{leancode}
\leanlabel{How the Lean proof follows the argument}
\texttt{lp\allowbreak{}Drift\_\allowbreak{}norm\_\allowbreak{}sq} rewrites \(\|\cdot\|^2\) as an inner product (\texttt{real\_\allowbreak{}inner\_\allowbreak{}self\_\allowbreak{}eq\_\allowbreak{}norm\_\allowbreak{}sq}) and the inner product as an integral (\texttt{L2.\allowbreak{}inner\_\allowbreak{}def}), then applies \texttt{integral\_\allowbreak{}congr\_\allowbreak{}ae} with the almost-everywhere identity \texttt{lp\allowbreak{}Drift\_\allowbreak{}coe}. Completeness of \(H\) is omitted from the hypotheses (\texttt{omit [Complete\allowbreak{}Space H] in}).
\end{leanbox}

\section{\texorpdfstring{A second-moment bound gives a drift bound}{A second-moment bound gives a drift bound}}


\textbf{Claim.} Under the hypotheses of the previous step, if \(C\in\mathbb R\) and \(\int\|g(X(\omega))\|^2\,P(d\omega)\le C\), then


\[\|\mathcal G(X)\|_{L^2}\le\sqrt C.\]


\textbf{Proof.} By the previous step \(\|\mathcal G(X)\|_{L^2}^2\le C\). For a real number \(a\) and any real \(s\), \(a^2\le s\) implies \(a\le\sqrt s\) (if \(s<0\) the hypothesis is impossible; otherwise \(a\le|a|=\sqrt{a^2}\le\sqrt s\)). \(\square\)




With \(g=\nabla U\) and \(C=Ld\), this converts the score moment bound into the hypothesis \(\|\mathcal G(Y_t)\|_{L^2}\le\sqrt{Ld}\) of the local-defect estimate.

\begin{leanbox}{Lean: \texttt{lpDrift\_\allowbreak{}norm\_\allowbreak{}le\_\allowbreak{}sqrt}}
\leansource{\texttt{lpDrift\_\allowbreak{}norm\_\allowbreak{}le\_\allowbreak{}sqrt} \quad UnadjustedLangevin/\allowbreak{}LpDrift.lean, lines 139–142}
\begin{leancode}
\leanline{0}{theorem\ lpDrift\_norm\_le\_sqrt\ (g\ :\ H\ →\ H)\ \{L\ :\ ℝ≥0\}\ (hLip\ :\ LipschitzWith\ L\ g)}
\leanline{4}{(X\ :\ Lp\ H\ 2\ P)\ (C\ :\ ℝ)\ (hmoment\ :\ (∫\ ω,\ ‖g\ (X\ ω)‖\ \textasciicircum{}\ 2\ ∂P)\ ≤\ C)\ :}
\leanline{4}{‖lpDrift\ P\ g\ hLip\ X‖\ ≤\ Real.sqrt\ C}
\end{leancode}
\leanlabel{How the Lean proof follows the argument}
\texttt{lp\allowbreak{}Drift\_\allowbreak{}norm\_\allowbreak{}le\_\allowbreak{}sqrt} is \texttt{Real.\allowbreak{}le\_\allowbreak{}sqrt\_\allowbreak{}of\_\allowbreak{}sq\_\allowbreak{}le} applied to \texttt{(lp\allowbreak{}Drift\_\allowbreak{}norm\_\allowbreak{}sq P g h\allowbreak{}Lip X).\allowbreak{}trans\_\allowbreak{}le hmoment}, i.e. to \(\|\mathcal G(X)\|^2=\int\|g(X)\|^2\le C\).
\end{leanbox}

\section{\texorpdfstring{The ULA update in \(L^2\)}{The ULA update in L²}}


\textbf{Definition.} Let \(H\) be a real Hilbert space, \(P\) a finite measure, \(U:H\to\mathbb R\) with \(L\)-Lipschitz gradient, \(\mathcal G\) the lift of \(\nabla U\), and \(h\in\mathbb R\). For \(X,Z\in L^2(P;H)\),


\[\mathsf T_h(X,Z)=X-h\,\mathcal G(X)+\sqrt{2h}\,Z\in L^2(P;H),\]


with \(\sqrt{2h}=0\) if \(h<0\). All three terms lie in \(L^2(P;H)\), which is a vector space.

\begin{leanbox}{Lean: \texttt{lpUlaUpdate}}
\leansource{\texttt{lpUlaUpdate} \quad UnadjustedLangevin/\allowbreak{}LpDrift.lean, lines 145–147}
\begin{leancode}
\leanline{0}{def\ lpUlaUpdate\ (U\ :\ H\ →\ ℝ)\ (h\ :\ ℝ)\ \{L\ :\ ℝ≥0\}}
\leanline{4}{(hLip\ :\ LipschitzWith\ L\ (gradient\ U))\ (X\ Z\ :\ Lp\ H\ 2\ P)\ :\ Lp\ H\ 2\ P\ :=}
\leanline{2}{X\ -\ h\ •\ lpDrift\ P\ (gradient\ U)\ hLip\ X\ +\ Real.sqrt\ (2\ *\ h)\ •\ Z}
\end{leancode}
\leanlabel{How the Lean proof follows the argument}
\texttt{lp\allowbreak{}Ula\allowbreak{}Update P U h h\allowbreak{}Lip X Z} is \texttt{X -\allowbreak{} h \ensuremath{\bullet} lp\allowbreak{}Drift P (gradient U) h\allowbreak{}Lip X + Real.\allowbreak{}sqrt (2 * h) \ensuremath{\bullet} Z}, computed with the vector-space operations of \texttt{Lp H 2 P}.
\end{leanbox}

\section{\texorpdfstring{The update acts pointwise}{The update acts pointwise}}


\textbf{Claim.} In the setting of the previous definition, for all \(X,Z\in L^2(P;H)\),


\[\mathsf T_h(X,Z)(\omega)=X(\omega)-h\nabla U(X(\omega))+\sqrt{2h}\,Z(\omega)\qquad\text{for }P\text{-almost every }\omega.\]


\textbf{Proof.} Outside the union of five null sets, the representatives of the sum, of the difference, of the two scalar multiples and of \(\mathcal G(X)\) (Step 3) are given by the corresponding pointwise formulas; substituting them gives the claim. \(\square\)




No assumption on the law of \(Z\) is involved.

\begin{leanbox}{Lean: \texttt{lpUlaUpdate\_\allowbreak{}coe}}
\leansource{\texttt{lpUlaUpdate\_\allowbreak{}coe} \quad UnadjustedLangevin/\allowbreak{}LpDrift.lean, lines 151–162}
\begin{leancode}
\leanline{0}{theorem\ lpUlaUpdate\_coe\ (U\ :\ H\ →\ ℝ)\ (h\ :\ ℝ)\ \{L\ :\ ℝ≥0\}}
\leanline{4}{(hLip\ :\ LipschitzWith\ L\ (gradient\ U))\ (X\ Z\ :\ Lp\ H\ 2\ P)\ :}
\leanline{4}{lpUlaUpdate\ P\ U\ h\ hLip\ X\ Z\ =ᵐ[P]\ fun\ ω\ ↦}
\leanline{6}{X\ ω\ -\ h\ •\ gradient\ U\ (X\ ω)\ +\ Real.sqrt\ (2\ *\ h)\ •\ Z\ ω}
\end{leancode}
\leanlabel{How the Lean proof follows the argument}
\texttt{lp\allowbreak{}Ula\allowbreak{}Update\_\allowbreak{}coe} unfolds the definition and uses \texttt{filter\_\allowbreak{}upwards} with \texttt{Lp.\allowbreak{}coe\allowbreak{}Fn\_\allowbreak{}add}, \texttt{Lp.\allowbreak{}coe\allowbreak{}Fn\_\allowbreak{}sub}, \texttt{Lp.\allowbreak{}coe\allowbreak{}Fn\_\allowbreak{}smul} (twice) and \texttt{lp\allowbreak{}Drift\_\allowbreak{}coe}; \texttt{simp only} substitutes the five identities.
\end{leanbox}

\section{\texorpdfstring{The ULA chain in \(L^2\)}{The ULA chain in L²}}


\textbf{Definition.} Let \(P\) be a finite measure, \(U:E_d\to\mathbb R\) with \(L\)-Lipschitz gradient, \(h\in\mathbb R\), \(X_0\in L^2(P;E_d)\), and \(Z_0,Z_1,\ldots\in L^2(P;E_d)\). The ULA chain in \(L^2\) is defined recursively by


\[X_0\ \text{given},\qquad X_{n+1}=\mathsf T_h(X_n,Z_n)\quad(n\ge0).\]


Each \(X_n\) belongs to \(L^2(P;E_d)\), because \(\mathsf T_h\) maps \(L^2\times L^2\) into \(L^2\).

\begin{leanbox}{Lean: \texttt{lpUlaChain}}
\leansource{\texttt{lpUlaChain} \quad UnadjustedLangevin/\allowbreak{}LpDrift.lean, lines 165–169}
\begin{leancode}
\leanline{0}{def\ lpUlaChain\ \{d\ :\ ℕ\}\ (U\ :\ E\ d\ →\ ℝ)\ (h\ :\ ℝ)\ \{L\ :\ ℝ≥0\}}
\leanline{4}{(hLip\ :\ LipschitzWith\ L\ (gradient\ U))\ (initial\ :\ Lp\ (E\ d)\ 2\ P)}
\leanline{4}{(noise\ :\ ℕ\ →\ Lp\ (E\ d)\ 2\ P)\ :\ ℕ\ →\ Lp\ (E\ d)\ 2\ P}
\leanline{2}{|\ 0\ =>\ initial}
\leanline{2}{|\ n\ +\ 1\ =>\ lpUlaUpdate\ P\ U\ h\ hLip\ (lpUlaChain\ U\ h\ hLip\ initial\ noise\ n)\ (noise\ n)}
\end{leancode}
\leanlabel{How the Lean proof follows the argument}
\texttt{lp\allowbreak{}Ula\allowbreak{}Chain P U h h\allowbreak{}Lip initial noise} is defined by structural recursion on \(n\): the clause for \texttt{0} returns \texttt{initial}, and the clause for \texttt{n + 1} returns \texttt{lp\allowbreak{}Ula\allowbreak{}Update P U h h\allowbreak{}Lip (lp\allowbreak{}Ula\allowbreak{}Chain U h h\allowbreak{}Lip initial noise n) (noise n)}. The noise sequence is \texttt{noise : \ensuremath{\mathbb N} \ensuremath{\to} Lp (E d) 2 P}.
\end{leanbox}

\section{\texorpdfstring{The initial iterate}{The initial iterate}}


\textbf{Claim.} The chain satisfies \(X_0=\) the given initial class.




\textbf{Proof.} This is the first clause of the recursive definition. \(\square\)

\begin{leanbox}{Lean: \texttt{lpUlaChain\_\allowbreak{}zero}}
\leansource{\texttt{lpUlaChain\_\allowbreak{}zero} \quad UnadjustedLangevin/\allowbreak{}LpDrift.lean, lines 171–173}
\begin{leancode}
\leanline{0}{@[simp]\ theorem\ lpUlaChain\_zero\ \{d\ :\ ℕ\}\ (U\ :\ E\ d\ →\ ℝ)\ (h\ :\ ℝ)\ \{L\ :\ ℝ≥0\}}
\leanline{4}{(hLip\ :\ LipschitzWith\ L\ (gradient\ U))\ (initial\ :\ Lp\ (E\ d)\ 2\ P)}
\leanline{4}{(noise\ :\ ℕ\ →\ Lp\ (E\ d)\ 2\ P)\ :\ lpUlaChain\ P\ U\ h\ hLip\ initial\ noise\ 0\ =\ initial}
\end{leancode}
\leanlabel{How the Lean proof follows the argument}
\texttt{lp\allowbreak{}Ula\allowbreak{}Chain\_\allowbreak{}zero} holds by \texttt{rfl} (definitional unfolding). The attribute \texttt{@[simp]} lets the simplifier use it automatically.
\end{leanbox}

\section{\texorpdfstring{The successor iterate}{The successor iterate}}


\textbf{Claim.} For every \(n\ge0\), \(X_{n+1}=\mathsf T_h(X_n,Z_n)\).




\textbf{Proof.} This is the second clause of the recursive definition. \(\square\)




The noise used in the step from \(n\) to \(n+1\) is \(Z_n\).

\begin{leanbox}{Lean: \texttt{lpUlaChain\_\allowbreak{}succ}}
\leansource{\texttt{lpUlaChain\_\allowbreak{}succ} \quad UnadjustedLangevin/\allowbreak{}LpDrift.lean, lines 175–179}
\begin{leancode}
\leanline{0}{@[simp]\ theorem\ lpUlaChain\_succ\ \{d\ :\ ℕ\}\ (U\ :\ E\ d\ →\ ℝ)\ (h\ :\ ℝ)\ \{L\ :\ ℝ≥0\}}
\leanline{4}{(hLip\ :\ LipschitzWith\ L\ (gradient\ U))\ (initial\ :\ Lp\ (E\ d)\ 2\ P)}
\leanline{4}{(noise\ :\ ℕ\ →\ Lp\ (E\ d)\ 2\ P)\ (n\ :\ ℕ)\ :}
\leanline{4}{lpUlaChain\ P\ U\ h\ hLip\ initial\ noise\ (n\ +\ 1)\ =}
\leanline{6}{lpUlaUpdate\ P\ U\ h\ hLip\ (lpUlaChain\ P\ U\ h\ hLip\ initial\ noise\ n)\ (noise\ n)}
\end{leancode}
\leanlabel{How the Lean proof follows the argument}
\texttt{lp\allowbreak{}Ula\allowbreak{}Chain\_\allowbreak{}succ} holds by \texttt{rfl} and is also a \texttt{@[simp]} lemma.
\end{leanbox}

\section{\texorpdfstring{The chain in \(L^2\) is the pointwise algorithm}{The chain in L² is the pointwise algorithm}}


\textbf{Claim.} In the setting of Step 13, let \(x_n(\omega)\) denote the pointwise ULA recursion driven by the representatives,


\[x_0(\omega)=X_0(\omega),\qquad x_{n+1}(\omega)=x_n(\omega)-h\nabla U(x_n(\omega))+\sqrt{2h}\,Z_n(\omega).\]


Then for every \(n\ge0\), \(X_n(\omega)=x_n(\omega)\) for \(P\)-almost every \(\omega\).




\textbf{Proof.} Induction on \(n\). For \(n=0\), both sides are the representative of the initial class, for every \(\omega\). Assume the claim for \(n\). By Step 12 applied to \(X_n\) and \(Z_n\), for almost every \(\omega\),


\[X_{n+1}(\omega)=X_n(\omega)-h\nabla U(X_n(\omega))+\sqrt{2h}\,Z_n(\omega).\]


Outside the union of this null set and the null set of the induction hypothesis, we may replace \(X_n(\omega)\) by \(x_n(\omega)\), and the right-hand side becomes \(x_{n+1}(\omega)\). \(\square\)




Hence the chain in \(L^2\) is the ordinary ULA algorithm applied to random inputs; in particular \(X_n\) has the same law as \(x_n\).

\begin{leanbox}{Lean: \texttt{lpUlaChain\_\allowbreak{}coe}}
\leansource{\texttt{lpUlaChain\_\allowbreak{}coe} \quad UnadjustedLangevin/\allowbreak{}LpDrift.lean, lines 183–193}
\begin{leancode}
\leanline{0}{theorem\ lpUlaChain\_coe\ \{d\ :\ ℕ\}\ (U\ :\ E\ d\ →\ ℝ)\ (h\ :\ ℝ)\ \{L\ :\ ℝ≥0\}}
\leanline{4}{(hLip\ :\ LipschitzWith\ L\ (gradient\ U))\ (initial\ :\ Lp\ (E\ d)\ 2\ P)}
\leanline{4}{(noise\ :\ ℕ\ →\ Lp\ (E\ d)\ 2\ P)\ (n\ :\ ℕ)\ :}
\leanline{4}{lpUlaChain\ P\ U\ h\ hLip\ initial\ noise\ n\ =ᵐ[P]\ fun\ ω\ ↦}
\leanline{6}{ulaChain\ U\ h\ (initial\ ω)\ (fun\ k\ ↦\ noise\ k\ ω)\ n}
\end{leancode}
\leanlabel{How the Lean proof follows the argument}
\texttt{lp\allowbreak{}Ula\allowbreak{}Chain\_\allowbreak{}coe} states the almost-everywhere equality with \texttt{ula\allowbreak{}Chain U h (initial \ensuremath{\omega}) (fun k \ensuremath{\mapsto} noise k \ensuremath{\omega}) n}, the pointwise chain defined in the section on the noise-driven iterates. The proof is \texttt{induction n}: the \texttt{zero} case is \texttt{Filter.\allowbreak{}Eventually.\allowbreak{}of\_\allowbreak{}forall fun \_\allowbreak{} \ensuremath{\mapsto} rfl}; in the \texttt{succ} case, \texttt{filter\_\allowbreak{}upwards} combines \texttt{lp\allowbreak{}Ula\allowbreak{}Update\_\allowbreak{}coe} for the \(n\)-th iterate with the induction hypothesis \texttt{ih}, and \texttt{simpa} unfolds both recursions (\texttt{lp\allowbreak{}Ula\allowbreak{}Chain\_\allowbreak{}succ}, \texttt{ula\allowbreak{}Chain\_\allowbreak{}succ}, \texttt{ula\allowbreak{}Update}) and substitutes \texttt{hprev}.
\end{leanbox}

\end{document}
